\section{Introduction}

Malaria caused by \textit{Plasmodium falciparum} remains a leading cause of morbidity and mortality in tropical regions, and genomic surveillance has emerged as a central tool for monitoring transmission, tracking drug resistance, and guiding elimination efforts \citep{who2022worldmalaria,malariagen2021pf6}. A foundational primitive of this surveillance is genetic relatedness: the probability that two parasite genomes share recent common ancestry at a given locus. Identity-by-descent (IBD) segments, contiguous stretches of chromosome inherited from a recent common ancestor, summarize this relatedness at fine temporal resolution and have become a workhorse for inferring effective population size ($N_e$), recent demography, and fine-scale population structure in \textit{P. falciparum} and other recombining pathogens \citep{schaffner2018hmmibd,henden2018isorelate,taylor2019ibd}.

IBD-based inferences rest on an underlying coalescent or Wright--Fisher model in which the distribution of segment lengths maps directly to the time to most recent common ancestor (TMRCA) and in which pairwise total IBD reflects genome-wide shared ancestry. The flagship $N_e$ estimator in this family, \texttt{IBDNe}, uses the empirical distribution of segment lengths to reconstruct $N_e$ trajectories over the past dozens of generations and has been widely applied to human and parasite populations \citep{browning2015ibdne,daniels2015modeling}. Parallel lines of work have repurposed IBD for community detection on relatedness networks \citep{rosvall2008maps}, for selection scans via positional enrichment statistics such as $X_{iR,s}$ \citep{henden2018xirs}, and for connecting population genetic structure to epidemiologic transmission \citep{chang2019mapping,wong2017genetic}.

The critical assumption shared by these methods is that the fraction of the genome shaped by natural selection is small relative to the neutral background, so that the IBD distribution primarily reflects demography. In \textit{P. falciparum}, however, this assumption is routinely violated. Parasites in Southeast Asia and, increasingly, in Africa, carry multiple simultaneous selective sweeps driven by antimalarial drugs: the \textit{kelch13} locus for artemisinin resistance \citep{ariey2014kelch13,ashley2014tracking}, \textit{pfcrt} for chloroquine and piperaquine \citep{fidock2000pfcrt,amato2017piperaquine}, and a genetic background of accessory loci including \textit{arps10}, \textit{pph}, \textit{ferredoxin}, and \textit{pfmdr1} \citep{miotto2015genetic}. Selective sweeps elongate the IBD segments that traverse the selected locus, inflate pairwise total IBD, and concentrate a disproportionate share of shared ancestry in a few chromosomal regions. Whether and how this distortion propagates into $N_e$ estimates and network-based structure inference has not been characterized, and there is no established correction protocol for IBD-based analyses in strongly selected populations.

A related literature on LD-based $N_e$ estimation has arrived at mixed conclusions. Classical theory predicts that selection shifts the LD decay curve and therefore biases $N_e$, yet \citet{novo2022ld} recently reported that LD-based estimators are virtually robust to natural selection across a broad range of parameters. Their simulations, however, used weak selection coefficients, randomly distributed selected loci, and long time scales, which capture a human-like regime rather than the concentrated, recent, strong selection typical of \textit{P. falciparum} under drug pressure. This apparent robustness therefore does not imply that IBD-based estimators, which depend on segment-length tails in a stronger way, inherit the same property.

In this work we directly quantify the effect of strong recent positive selection on IBD-based inferences in \textit{P. falciparum} and evaluate a simple correction based on IBD peak removal. Our contributions are fourfold. First, we assemble whole-genome sequencing data from 640 new and 4,026 publicly available clinical isolates, yielding an analyzable Southeast Asian (SEA) dataset of 2,055 parasites distributed across 14 years and 18 provinces in Cambodia, Thailand, Laos, and Vietnam, and a validation West African (WAF) dataset from \textit{P. falciparum} catalogue \citep{malariagen2021pf6,manske2012analysis}. Second, we introduce \texttt{tskibd}, a true-IBD algorithm that operates on simulated genealogical trees in tree-sequence format \citep{kelleher2018efficient,haller2019treeseq,baumdicker2022msprime} and therefore decouples questions about IBD quality from questions about selection bias. Third, we use single-population and multi-deme stepping-stone simulations \citep{haller2019slim,kelleher2016msprime} to show that a representative selection scenario ($s = 0.3$, selection beginning 80 generations ago, single-origin favoured allele at 33.3~cM on each chromosome) distorts segment-length, pairwise total, and chromosomal positional IBD distributions, underestimates recent $N_e$ via \texttt{IBDNe}, and blurs InfoMap-detected community structure. Fourth, we develop and evaluate an IBD peak removal strategy, showing that the correction recovers accurate $N_e$ and structure in low-background-relatedness populations (WAF) but has little effect and risks over-correction in high-background-relatedness populations (SEA) \citep{daniels2015modeling,verity2020diversity}.

\section{Related Work}

\paragraph{IBD inference in \textit{P. falciparum}.} Two hidden Markov model callers dominate haploid parasite IBD inference: \texttt{hmmIBD} \citep{schaffner2018hmmibd} and \texttt{isoRelate} \citep{henden2018isorelate}. Both emit a binary per-site IBD state along pairs of haplotypes and summarise relatedness via the fraction of sites called IBD. Segment-level properties, although reported, are not routinely validated. \citet{taylor2019ibd} present a Bayesian framework that integrates relatedness uncertainty for downstream estimators.

\paragraph{Demographic estimators.} \texttt{IBDNe} \citep{browning2015ibdne} uses the empirical distribution of IBD segment lengths to estimate $N_e$ over the recent past, and has been deployed on \textit{P. falciparum} data from high- and low-transmission settings \citep{daniels2015modeling,wong2017genetic}. LD-based alternatives place different weight on recombination-distance bins and have been argued to be robust to selection in simulated human-like settings \citep{novo2022ld}; this claim is the closest analogue to ours but reaches a different conclusion under a different simulation regime, which we discuss in the main text.

\paragraph{Selection detection with IBD.} IBD-based selection scans exploit positional enrichment of segments near favoured alleles. The $X_{iR,s}$ statistic of \citet{henden2018xirs} summarises this enrichment relative to a null expectation and motivates our peak-identification procedure, which applies a chromosome-wise threshold and a boundary-extension heuristic to flag sweep regions for removal.

\paragraph{Tree-sequence simulation.} Our true-IBD pipeline builds on modern tree-sequence tooling: \texttt{msprime} \citep{kelleher2016msprime,baumdicker2022msprime} for coalescent backgrounds, \texttt{SLiM} \citep{haller2019slim} with tree-sequence recording \citep{haller2019treeseq,kelleher2018efficient} for selection and structured demographies, and \texttt{tskit} data structures throughout. Recording the genealogy in simulation lets us bypass the confounding between IBD caller accuracy and the selection effects we wish to isolate.

\paragraph{Network community detection.} We use \texttt{InfoMap} \citep{rosvall2008maps} for community detection on IBD-weighted networks, evaluated against ground-truth labels via modularity \citep{newman2004modularity} and the adjusted Rand index. Within-community and inter-community IBD ratios characterise the strength of population boundaries before and after peak correction.

\paragraph{Empirical context.} Genomic surveillance in \textit{P. falciparum} has documented the spread of artemisinin resistance across the Greater Mekong Subregion \citep{ashley2014tracking,miotto2015genetic} and the sustained high relatedness of SEA parasite populations as transmission declined \citep{amato2017piperaquine}. West Africa, by contrast, has higher $N_e$ and lower background relatedness, even as Kelch13 variants have emerged locally \citep{verity2020diversity}. The contrast between these settings is central to our finding that background relatedness is the key modifier of selection-induced IBD bias and the utility of correction.
